\documentclass[10pt,a4paper]{amsart}
\usepackage[margin=19mm]{geometry}
\usepackage{amsmath,amssymb,mathtools}
\usepackage[scr=rsfs,cal=euler]{mathalfa}
\usepackage{xfrac}
\usepackage[unicode,hidelinks,bookmarks=false]{hyperref}
\newcommand{\ie}{i.e.\ }
\newcommand{\cf}{cf.\ }
\newcommand{\resp}{respectively\ }
\newcommand{\Subsection}{\subsection}
\providecommand{\nobreakdash}{\nobreak\hbox{-}}
\setlength{\emergencystretch}{2em}
\multlinegap0pt
\usepackage{amssymb}
\usepackage[shortlabels]{enumitem}
\usepackage{xcolor}
\newtheorem{lemma}{Lemma}
\title{\textbf{Boundary Geometry and Surjective Linear Isometries of Weighted Hardy Spaces}}
\author{Ren-Yu Chen, Song-Ying Li, Sujuan Long and Jie Luo$^{*}$}
\date{Source: arXiv:2609.04612v1 (CC BY 4.0)}
\begin{document}
\maketitle

\noindent\emph{Source and licence.} \textbf{Boundary Geometry and Surjective Linear Isometries of Weighted Hardy Spaces}. Version arXiv:2609.04612v1 (CC BY 4.0), distributed by arXiv under the Creative Commons Attribution 4.0 International licence, http://creativecommons.org/licenses/by/4.0/, as recorded in the arXiv licence metadata for this exact version and shown on the abstract page https://arxiv.org/abs/2609.04612v1. Full licence notice: \texttt{licenses/arxiv-2609.04612-CC-BY-4.0.txt}.\par\smallskip

\noindent\textit{Editorial scope.} Counted unit: the article's lemma 2.1 (source0.tex lines 414--428). The article's complete original proof of it is delivered with it (source0.tex lines 430--522, one \texttt{proof} environment of 93 lines). No mathematics is written, repaired or re-derived: the statement and the proof are the article's own lines, in the article's own notation, and no retained line is substituted or re-flowed.  No further source-local context is needed: every cross-reference of every retained line resolves inside the two delivered spans.  Nothing was omitted from any retained span, so the package carries no omission record.  No retained line contains a citation, so the wrapper carries no bibliography and no bibliography entry.  No retained line contains a figure environment, an image-inclusion command or drawing code, so no image is delivered and the package declares no asset dependency.  Editorial formatting only, in addition: an AMS \texttt{amsart} preamble that restates the article's own declarations and macros used by the retained lines, namely the environments \texttt{lemma} on separate counters, together with the standard packages those lines need.  The retained environments print in this document as Lemma 1 in order of appearance, whereas the article prints the same items with the numbering of its own full document.  The complete native source of the article as distributed in the arXiv e-print for this exact version is bundled unchanged as \texttt{sources/209376/source0.tex}.
\begin{lemma}
\label{2.1} Let $D\subset\mathbb{C}^{n}$ be a bounded domain with $C^{2}$
boundary. Then
%there exists a defining function $\rho$ of $D$ such that for
every $f\in H(D)$, the quantity
\[
\int_{\partial D_{\varepsilon}}|f(\xi)|^{p}
%\frac{\partial\rho}{\partial\nu}(\xi)\,
\omega(\xi) d\sigma_{\varepsilon}(\xi)
\]
is decreasing for sufficiently small $0<\varepsilon<\hbox{dist}(z_{0},
\partial D)/2$, where $z_{0}\in D$, $\omega(z)={\frac{\partial G(z, z_{0}%
)}{\partial\nu(z)}}$ and $G(z, z_{0})$ is the Green function of the Laplacian
$\Delta$.
\end{lemma}

\begin{proof}
For $z_{0}\in D$, let $w(z)=G(z, z_{0})$ be the solution of
\[%
\begin{cases}
\Delta w(z)=\delta(z-z_{0}) & \text{in }D,\\
w(\zeta)=0 & \text{on }\partial D.
\end{cases}
\]
By the maximum principle or Hopf lemma,
\[
\omega(z)=\frac{\partial w}{\partial\nu}>0\qquad\text{on }\partial D.
\]
Let $0<\varepsilon_{1}<\varepsilon_{2}<\hbox{dist}(z_{0}, \partial D)$. Then
by the Green's formula
\begin{align*}
&  \int_{\partial D_{\varepsilon_{1}}}|f(\xi)|^{p}\frac{\partial w}%
{\partial\nu_{\varepsilon_{1}}}(\xi)\,d\sigma_{\varepsilon_{1}}(\xi
)-\int_{\partial D_{\varepsilon_{2}}}|f(\xi)|^{p}\frac{\partial w}{\partial
\nu_{\varepsilon_{2}}}(\xi)\,d\sigma_{\varepsilon_{2}}(\xi)\\
&  \qquad-\left(  \int_{\partial D_{\varepsilon_{1}}}\frac{\partial|f|^{p}%
}{\partial\nu_{\varepsilon_{1}}}(\xi)\,w(\xi)\,d\sigma_{\varepsilon_{1}}%
(\xi)-\int_{\partial D_{\varepsilon_{2}}}\frac{\partial|f|^{p}}{\partial
\nu_{\varepsilon_{2}}}(\xi)\,w(\xi)\,d\sigma_{\varepsilon_{2}}(\xi)\right) \\
&  =\int_{D_{\varepsilon_{1}}\setminus D_{\varepsilon_{2}}}|f|^{p}(z)\Delta
w(z)\,dv(z)-\int_{D_{\varepsilon_{1}}\setminus D_{\varepsilon_{2}}}%
\Delta(|f|^{p})(z)\,w(z)\,dv(z).
\end{align*}
Since $w$ is harmonic on $D_{\varepsilon_{1}}\setminus\overline{D_{\varepsilon
_{2}}},$ we have $\Delta w=0$ there. Hence
\begin{align*}
&  \int_{\partial D_{\varepsilon_{1}}}|f|^{p}\frac{\partial w}{\partial
\nu_{\varepsilon_{1}}}\,d\sigma_{\varepsilon_{1}}-\int_{\partial
D_{\varepsilon_{2}}}|f|^{p}\frac{\partial w}{\partial\nu_{\varepsilon_{2}}%
}\,d\sigma_{\varepsilon_{2}}\\
&  =-\int_{D_{\varepsilon_{1}}\setminus D_{\varepsilon_{2}}}\Delta
(|f|^{p})(z)\,w(z)\,dv(z)+\Big(\int_{\partial D_{\varepsilon_{1}}}%
\frac{\partial|f|^{p}}{\partial\nu_{\varepsilon_{1}}}(\xi)\,w(\xi
)\,d\sigma_{\varepsilon_{1}}(\xi)-\int_{\partial D_{\varepsilon_{2}}}%
\frac{\partial|f|^{p}}{\partial\nu_{\varepsilon_{2}}}(\xi)\,w(\xi
)\,d\sigma_{\varepsilon_{2}}(\xi)\Big).
\end{align*}
Since $w=-\varepsilon_{j}\ \text{on }\partial D_{\varepsilon_{j}},$ it
follows that
\begin{align*}
&  \int_{\partial D_{\varepsilon_{1}}}|f|^{p}\frac{\partial w}{\partial
\nu_{\varepsilon_{1}}}\,d\sigma_{\varepsilon_{1}}-\int_{\partial
D_{\varepsilon_{2}}}|f|^{p}\frac{\partial w}{\partial\nu_{\varepsilon_{2}}%
}\,d\sigma_{\varepsilon_{2}}\\
&  =-\int_{D_{\varepsilon_{1}}\setminus D_{\varepsilon_{2}}}\Delta
(|f|^{p})(z)\,w(z)\,dv(z)-\varepsilon_{1}\int_{\partial D_{\varepsilon_{1}}%
}\frac{\partial|f|^{p}}{\partial\nu_{\varepsilon_{1}}}\,d\sigma_{\varepsilon
_{1}}+\varepsilon_{2}\int_{\partial D_{\varepsilon_{2}}}\frac{\partial|f|^{p}%
}{\partial\nu_{\varepsilon_{2}}}\,d\sigma_{\varepsilon_{2}}.
\end{align*}
Applying Green's identity again,
\[
\int_{\partial D_{\varepsilon_{j}}}\frac{\partial|f|^{p}}{\partial
\nu_{\varepsilon_{j}}}\,d\sigma_{\epsilon_j}=\int_{D_{\varepsilon_{j}}}\Delta
(|f|^{p})(z)\,dv(z),\qquad j=1,2.
\]
Therefore
\begin{align*}
\lefteqn{\int_{\partial D_{\varepsilon_{1}}}\left\vert f\right\vert ^{p}%
\frac{\partial w}{\partial\nu_{\varepsilon_{1}}}d\sigma_{\varepsilon_{1}}%
-\int_{\partial D_{\varepsilon_{2}}}\left\vert f\right\vert ^{p}\frac{\partial
w}{\partial\nu_{\varepsilon_{2}}}d\sigma_{\varepsilon_{2}}}\\
&  =-\int_{D_{\varepsilon_{1}}\backslash D_{\varepsilon_{2}}}\Delta\left(
\left\vert f\right\vert ^{p}\right)  wdv-\varepsilon_{1}\int_{D_{\varepsilon
_{1}}}\Delta\left\vert f\right\vert ^{p}dv+\varepsilon_{2}\int_{D_{\varepsilon
_{2}}}\Delta\left\vert f\right\vert ^{p}dv\\
&  =-\int_{D_{\varepsilon_{1}}\backslash D_{\varepsilon_{2}}}\Delta\left(
\left\vert f\right\vert ^{p}\right)  \left(  w+\varepsilon_{1}\right)
dv+\left(  \varepsilon_{2}-\varepsilon_{1}\right)  \int_{D_{\varepsilon_{2}}%
}\Delta\left\vert f\right\vert ^{p}dv.
\end{align*}
Since $|f|^{p}$ is subharmonic and
\[
-\varepsilon_{2}\leq w(z)\leq-\varepsilon_{1}\quad\mbox{on}\quad
D_{\varepsilon_{1}}\backslash D_{\varepsilon_{2}},
\]
we have
\[
w(z)+\varepsilon_{1}\leq0\quad\hbox{and}\quad\Delta(|f|^{p})\geq0.
\]
Consequently,
\[
\int_{\partial D_{\varepsilon_{1}}}|f|^{p}\frac{\partial w}{\partial
\nu_{\varepsilon_{1}}}\,d\sigma_{\varepsilon_{1}}-\int_{\partial
D_{\varepsilon_{2}}}|f|^{p}\frac{\partial w}{\partial\nu_{\varepsilon_{2}}%
}\,d\sigma_{\varepsilon_{2}}\geq0.
\]
This proves the monotonicity.
\end{proof}

\end{document}
